\section{Introduction}
\label{sec:intro}

A customer support desk \af{should be the next ticket to attend}must decide which next ticket should be attended to. High severity issues such as payment fraud require immediate intervention under strict response windows, \af{no need to mention in minutes} whereas routine balance inquiries tolerate delays of several hours. When customer requests arrive faster than staff can process them, a backlog forms. In this operating environment, the dispatch order in which tickets are assigned to agents directly determines whether service level agreements (SLAs) are upheld or violated.

Published work on bank complaint handling has focused on predicting static categories and evaluating held-out classification performance \cite{arya2026bank}. A predicted label alone, however, does not specify a dispatch order. One possible policy serves tickets First Come First Served (FCFS); another groups them by predicted priority and uses arrival order within each group. Deshmukh and Shiravale \cite{deshmukh2020priority}, for example, prioritize complaints by sentiment intensity, break equal-intensity ties by arrival order, and introduce a waiting-time threshold for lower-intensity complaints.

Converting classifier outputs into an operational queue raises three challenges: prolonged waits under static priority, loss of predictive uncertainty, and variation across languages and scripts. First, static priority allows newly arriving urgent inquiries to overtake older routine requests. This can leave low-priority customers with inadequate service levels; delay-dependent and accumulating-priority disciplines address this by incorporating elapsed waiting time \cite{kleinrock1964delay,stanford2014accumulating}. Second, discrete argmax binning places a confident High-priority prediction (0.99 posterior probability) and an uncertain 0.51 versus 0.49 near tie into the same tier. In a two-type queue with imperfect class information and linear waiting costs, Argon and Ziya \cite{argon2009priority} show that ranking customers by their posterior signal outperforms any finite-class priority policy.

Third, multilingual banking messages span English, Sinhala, and Tamil, alongside informal Romanized scripts (Singlish and Tanglish). Sinhala-English code-mixed language identification and Sinhala script-sensitivity studies document challenges associated with these script variants \cite{smith2020codemixed,rajapakse2026script}. BANKING77 provides English queries labeled for 77 banking intents \cite{casanueva2020banking77}; it does not provide parallel Sinhala, Tamil, or Romanized versions with joint sentiment and priority annotations. We therefore extend this corpus to support multilingual classification and dispatch research.

Under this operational setting, we examine two research questions:

\textbf{RQ1} How accurately do fine-tuned multilingual pre-trained models (sentence encoders, masked encoders and small decoders) classify the intent, sentiment and priority of banking tickets in Sinhala, Tamil and their romanized forms, compared to a TF–IDF baseline? \af{not clear. is it empirically evaluate exisistng multilingual pre-trained models accuracy for .....? Dont bring class imbalance here. RQ should be specific and measureable}

\textbf{RQ2} How can continuous classifier posteriors for ticket severity, sentiment, and intent criticality be combined into a single urgency score? \af{dont think the waiting time is a consideration. It is external factor. Frame problem as to dynamically decide urgency based on .....what?}

Our contributions are twofold. First, a \textbf{multilingual model benchmark study} spanning classical TF--IDF models, sentence encoders, masked encoders, and instruction-tuned decoders over five language variants\af{what are tracks? Dataset?}. Second, a \textbf{urgency scoring method} that combines the output distributions of the separate intent and priority classifiers to a single priority distribution via logarithmic opinion pooling.\af{what is this?}
